\chapter{数学速查：为模型计算服务}
\label{ch:math}

本附录保留理解本书公式最常用的数学工具：形状、内积、矩阵运算、导数、概率与统计量。
每个公式都应先确认变量含义与适用条件。无需在开始学习机器学习之前逐页背诵。

\begin{keyidea}[先检查形状，再检查数值]
矩阵相乘看内部维度，梯度应与参数具有相同形状，概率应满足非负与归一化。
这三类检查能较早发现许多实现错误。
\end{keyidea}

\section{向量、矩阵与乘法形状}
\label{sec:math-matrix}
列向量$\vct x\in\R^d$有$d$个分量；矩阵$A\in\R^{m\times n}$有$m$行$n$列，
元素记为$a_{ij}$，转置满足$(A^\top)_{ji}=a_{ij}$。
矩阵积$C=AB$的元素为$c_{ij}=\sum_{k=1}^n a_{ik}b_{kj}$。

\begin{figure}[H]
\centering\begin{tikzpicture}[node distance=8mm]
\node[box,text width=30mm] (a) {$A_{2\times3}$\\2行，3个分量};
\node[box,text width=30mm,right=of a] (b) {$B_{3\times2}$\\3个分量，2列};
\node[box,text width=30mm,right=12mm of b] (c) {$C_{2\times2}$\\2行，2列};
\node[font=\Large] at ($(a.east)!0.5!(b.west)$) {$\times$};
\draw[flow] (b)--(c);
\node[lab,below=10mm of b] {内部维度相同才能相乘；结果保留外部维度};
\end{tikzpicture}

\caption{例子$A\in\R^{2\times3}$、$B\in\R^{3\times2}$，相乘后为$2\times2$。}
\end{figure}
\readfigure{结果每个位置都是左矩阵一行与右矩阵一列的内积。一般$AB\ne BA$，
甚至可能只有一个方向的积有定义。维度合适时，乘法满足结合律。}
\[
 A=\begin{bmatrix}1&2&0\\0&1&3\end{bmatrix},\quad
 B=\begin{bmatrix}2&1\\1&2\\3&0\end{bmatrix},\quad
 AB=\begin{bmatrix}4&5\\10&2\end{bmatrix}.
\]
例如$(AB)_{12}=1\times1+2\times2+0\times0=5$。
逐元素乘法$A\odot B$要求相应形状匹配，不是上述矩阵乘法。

\quick{批量线性预测常写作$X\vct w+b\vct1$：$X$为$n\times d$，$\vct w$为$d\times1$，结果为$n\times1$。}

\clearpage
\section{内积、范数、秩与行列式}
\label{sec:math-norms}
内积$\vct x^\top\vct z=\sum_jx_jz_j$衡量方向上的关系。
常用范数为
\[
 \lVert\vct x\rVert_1=\sum_j|x_j|,\quad
 \lVert\vct x\rVert_2=\sqrt{\sum_jx_j^2},\quad
 \lVert A\rVert_F=\sqrt{\sum_{i,j}a_{ij}^2}.
\]
例如$(3,-4)$的L1范数为7，L2范数为5；范数总是非负。
矩阵的秩（Rank）是独立列或独立行的最大数目，不等于行数、列数的简单乘积。
向量的线性张成空间（Span）是它们所有线性组合构成的集合，\textbf{不是夹角}。
例如$(1,0)$和$(2,0)$只能张成横轴；$(1,0)$和$(0,1)$则张成整个$\R^2$。

\begin{figure}[H]
\centering\begin{tikzpicture}[x=18mm,y=18mm]
\draw[->,Quiet] (-.2,0)--(3.5,0) node[right,lab] {$x_1$};
\draw[->,Quiet] (0,-.2)--(0,2.5) node[above,lab] {$x_2$};
\filldraw[fill=Pale,draw=Teal] (0,0)--(2,0)--(3,2)--(1,2)--cycle;
\draw[flow,very thick] (0,0)--(2,0) node[midway,below,lab] {$\vct a=(2,0)$};
\draw[flow,very thick] (0,0)--(1,2) node[midway,left,lab] {$\vct b=(1,2)$};
\node[lab] at (2,1) {面积4};
\end{tikzpicture}

\caption{二维列向量构成的平行四边形，面积为行列式绝对值。}
\end{figure}
\readfigure{$A=[\vct a\ \vct b]=\left[\begin{smallmatrix}2&1\\0&2\end{smallmatrix}\right]$，
$\det A=2\times2-1\times0=4$，面积为4。行列式本身有符号，交换两列会变号。}

$n\times n$方阵的行列式是$n$个特征值的乘积，\textbf{不是所有元素的乘积}。
$\det A\ne0$等价于方阵可逆。非方阵没有此处定义的行列式。
实际求解线性方程常采用稳定分解方法，不必先显式计算$A^{-1}$。

\clearpage
\section{特征值与奇异值：方向、缩放与降维}
\label{sec:math-spectral}
若非零向量$\vct v$满足$A\vct v=\lambda\vct v$，则$\lambda$是特征值，
$\vct v$是特征向量。实对称矩阵可写成$A=Q\Lambda Q^\top$，$Q$的列为正交单位向量。

\begin{center}
\begin{tabularx}{\textwidth}{P{32mm}Y}
\toprule
工具 & 用途与条件\\
\midrule
正半定 & 对称矩阵满足$\vct z^\top A\vct z\geq0$；等价于特征值非负。\\
奇异值分解 SVD & 任意实$m\times n$矩阵可写成$A=U\Sigma V^\top$，奇异值非负。\\
低秩近似 & 只保留最大的若干奇异值和相应方向，以压缩矩阵。\\
PCA的线性代数 & 对中心化的数据矩阵作SVD，或对相应协方差矩阵求特征方向。\\
\bottomrule
\end{tabularx}
\end{center}

例如$A=\operatorname{diag}(3,1)$把两个坐标方向分别缩放为3倍、1倍，
奇异值为3、1。保留第一方向得到$A_1=\operatorname{diag}(3,0)$，
误差$\lVert A-A_1\rVert_F=1$。

\begin{pitfall}[方向不是自动“有用”的标签信息]
PCA寻找较大方差方向，不直接使用目标标签，也不保证保留最强预测信号。
用于监督流程时，中心化和降维都应在训练内拟合。
一般矩阵不一定有实数特征值；SVD与特征分解不是同一概念。
\end{pitfall}

\clearpage
\section{导数、梯度与链式法则}
\label{sec:math-gradient}
一元导数$f'(x)$描述局部变化率；多元标量函数的梯度由各偏导组成：
\[
\nabla_{\vct w}J=(\partial J/\partial w_1,\ldots,\partial J/\partial w_d)^\top.
\]
在可微条件下，$J(\vct w+\Delta)\approx J(\vct w)+\nabla J(\vct w)^\top\Delta$。

\begin{center}
\begin{tabular}{ll}
\toprule
函数 & 导数或梯度\\
\midrule
$x^p$ & $p x^{p-1}$（在有定义且可微处）\\
$e^x$ & $e^x$\\
$\log x$ & $1/x$，$x>0$\\
$\sigma(x)=1/(1+e^{-x})$ & $\sigma(x)(1-\sigma(x))$\\
$\frac12\lVert\vct w\rVert_2^2$ & $\vct w$\\
$\frac1{2n}\lVert X\vct w-\vct y\rVert_2^2$ & $\frac1nX^\top(X\vct w-\vct y)$\\
\bottomrule
\end{tabular}
\end{center}

复合函数$z=f(g(x))$的导数为$f'(g(x))g'(x)$。
多条路径依赖同一变量时，各路径贡献相加，这正是计算图中反向传播的基础。

\subsection{一个形状与数值检查}
$X=\left[\begin{smallmatrix}1&2\\3&4\end{smallmatrix}\right]$、
$\vct w=(0.5,-0.25)^\top$、$\vct y=(1,0)^\top$，则残差为$(-1,0.5)^\top$，
$J=0.3125$、$\nabla J=(0.25,0)^\top$。梯度与$\vct w$均为两维列向量。

\begin{pitfall}[不可微与零梯度]
$|x|$在0处不可微，可以使用次梯度；ReLU在0处也需实现约定。
梯度为0不一定是极小值。神经网络反向传播要先正确设置损失与计算依赖。
\end{pitfall}

\clearpage
\section{条件概率与贝叶斯公式}
\label{sec:math-probability}
对事件$A,B$，$P(A\cup B)=P(A)+P(B)-P(A\cap B)$。
当$P(B)>0$时，$P(A\mid B)=P(A\cap B)/P(B)$。
独立要求$P(A\cap B)=P(A)P(B)$；互斥事件并不因此独立。

设$B_1,\ldots,B_K$构成互斥且完备的划分，$P(B_k)>0$，则
\[
 P(A)=\sum_kP(A\mid B_k)P(B_k),\qquad
 P(B_j\mid A)=\frac{P(A\mid B_j)P(B_j)}{P(A)}.
\]
后一式还要求$P(A)>0$。先验$P(B_j)$、似然$P(A\mid B_j)$与后验$P(B_j\mid A)$的方向不同。

\begin{center}
\begin{tabularx}{\textwidth}{P{43mm}Y}
\toprule
某检测的已知概率 & 数值\\
\midrule
患病先验 & $P(D)=0.01$\\
患病者检出率 & $P(+\mid D)=0.90$\\
未患病者假阳性率 & $P(+\mid\bar D)=0.05$\\
\bottomrule
\end{tabularx}
\end{center}
\[
 P(D\mid+)=\frac{0.90\times0.01}{0.90\times0.01+0.05\times0.99}
 =\frac2{13}\approx0.154.
\]
检出率90\%并不意味着阳性者有90\%患病；先验与假阳性共同影响后验。

\quick{遇到条件概率，先读清楚竖线两侧，再决定能否交换方向。}

\clearpage
\section{随机变量、期望、方差与协方差}
\label{sec:math-random}
随机变量（Random Variable）把试验结果映射到数值。
例如两次抛硬币的正面个数，把四种等可能结果映射到0、1、2。

\begin{figure}[H]
\centering\begin{tikzpicture}[x=1mm,y=1mm]
\draw[Teal,fill=Pale] (0,0) rectangle (39,48);
\node[lab] at (19.5,53) {试验结果};
\foreach \y/\out in {39/反反,29/正反,19/反正,9/正正}{\node at (19,\y) {\out};}
\draw[->,Quiet] (67,0)--(130,0) node[right,lab] {$X$};
\foreach \x/\v in {72/0,96/1,120/2}{\draw[Quiet] (\x,0)--(\x,3);\node[lab,below] at (\x,-1) {$\v$};}
\draw[flow] (39,39) to[bend left=15] (72,3);
\draw[flow] (39,29) to[bend left=15] (96,3);
\draw[flow] (39,19) to[bend left=10] (96,3);
\draw[flow] (39,9) to[bend left=10] (120,3);
\node[lab] at (94,49) {$X=$两次抛硬币的正面数};
\node[lab] at (94,39) {“正反”与“反正”都映射到1};
\end{tikzpicture}

\caption{四种试验结果映射到三个计数值；不同结果可以映射到相同数值。}
\end{figure}
\readfigure{公平且独立的两次抛掷中，$P(X=0)=1/4$、$P(X=1)=1/2$、$P(X=2)=1/4$，
所以$\mathbb E X=1$、$\operatorname{Var}(X)=1/2$。}

离散随机变量有$\mathbb E[g(X)]=\sum_x g(x)P(X=x)$；连续情形用密度积分。
若相应二阶矩存在，
\[
 \operatorname{Var}(X)=\mathbb E[(X-\mathbb E X)^2]=\mathbb E[X^2]-(\mathbb E X)^2,
\]
\[
 \operatorname{Cov}(X,Y)=\mathbb E[(X-\mathbb E X)(Y-\mathbb E Y)].
\]
期望满足线性性，不要求独立；独立且二阶矩有限时协方差为0，逆命题一般不成立。

\begin{pitfall}[零相关仍可能有确定依赖]
令$X$等可能取$-1,0,1$，$Y=X^2$，则协方差为0，
但$Y$完全由$X$决定。相关系数只反映相应的线性关系，且两变量方差应非零。
\end{pitfall}

\clearpage
\section{常用分布与样本统计量}
\label{sec:math-distributions}
概率质量函数直接给离散取值的概率；概率密度需要在区间上积分。
密度可以大于1，例如$U(0,0.5)$的密度为2，整个区间概率仍为1。

\begin{center}
\begin{tabularx}{\textwidth}{P{31mm}P{67mm}Y}
\toprule
分布 & 质量函数或密度 & 均值与方差\\
\midrule
Bernoulli$(p)$ & $p^x(1-p)^{1-x}$，$x\in\{0,1\}$ & $p$；$p(1-p)$\\
Binomial$(n,p)$ & $\binom nx p^x(1-p)^{n-x}$，$x=0,\ldots,n$ & $np$；$np(1-p)$\\
Uniform$(a,b)$ & $1/(b-a)$在$[a,b]$内，区间外0 & $(a+b)/2$；$(b-a)^2/12$\\
Normal$(\mu,\sigma^2)$ & $\exp[-(x-\mu)^2/(2\sigma^2)]/(\sigma\sqrt{2\pi})$ & $\mu$；$\sigma^2$\\
\bottomrule
\end{tabularx}
\end{center}
$0\leq p\leq1$，二项$n$为非负整数，均匀分布$a<b$，正态$\sigma>0$。
固定$n$次独立同分布Bernoulli试验的和才具有上述二项分布。

\subsection{总体矩与样本估计要区分}
对独立同分布观测$x_1,\ldots,x_n$，样本均值为$\bar x=\sum_i x_i/n$。
$n>1$时，$s^2=\sum_i(x_i-\bar x)^2/(n-1)$是总体方差的无偏估计（有限二阶矩）。
描述手头数据或某些标准化实现则使用分母$n$，应明确约定。
对$[1,2,3]$，均值为2，分母$n$的方差为$2/3$，分母$n-1$为1。

\begin{pitfall}[大样本结论仍有条件]
在独立同分布等适当条件下，大数定律说明均值趋近总体期望，
中心极限定理描述标准化均值的近似分布。它们不表示任意原始数据都会随着样本增多变成正态。
\end{pitfall}

\section{CDF：把“落在左侧的概率”累计起来}
\label{sec:math-cdf}
累积分布函数（Cumulative Distribution Function, CDF）定义为
\[
 F_X(t)=P(X\leq t),\quad t\in\R.
\]
$X$是随机变量，$t$是查询位置；不要把等式右侧写成与$t$无关的固定事件。
CDF单调不减、右连续，在负无穷和正无穷处的极限分别为0和1。
它适用于离散、连续以及混合分布。

\begin{figure}[htbp]
\centering\begin{tikzpicture}
\begin{axis}[name=disc,width=72mm,height=43mm,xmin=-.3,xmax=1.3,ymin=-.05,ymax=1.12,
 xlabel={$t$},ylabel={$F_X(t)$},xtick={0,1},ytick={0,.75,1},axis lines=left]
\addplot[Teal,thick] coordinates {(-.3,0) (0,0)};
\addplot[Teal,thick] coordinates {(0,.75) (1,.75)};
\addplot[Teal,thick] coordinates {(1,1) (1.3,1)};
\addplot[only marks,mark=*,Teal,mark size=2pt] coordinates {(0,.75) (1,1)};
\addplot[only marks,mark=o,mark options={fill=white},Teal,mark size=2pt] coordinates {(0,0) (1,.75)};
\end{axis}
\begin{axis}[at={(disc.east)},anchor=west,xshift=13mm,width=72mm,height=43mm,
 xmin=-.15,xmax=.7,ymin=-.05,ymax=1.12,xlabel={$t$},xtick={0,.5},ytick={0,1},axis lines=left]
\addplot[Teal,thick] coordinates {(-.15,0) (0,0) (.5,1) (.7,1)};
\addplot[only marks,mark=*,Teal,mark size=2pt] coordinates {(0,0) (.5,1)};
\end{axis}
\end{tikzpicture}

\caption{左：$X\sim\mathrm{Bernoulli}(1/4)$的跳跃CDF。右：$U(0,0.5)$的连续CDF。实心点表示端点取值。}
\par\smallskip
\readfigure{左图在0跳到$3/4$，在1跳到1；所以$P(X=0)=3/4$。
右图在$[0,0.5]$内为$F(t)=2t$，因此$P(0.1<X\leq0.4)=0.8-0.2=0.6$。}
\end{figure}

\begin{center}
\begin{tabularx}{\textwidth}{P{30mm}Y}
\toprule
对象 & 与概率的关系\\
\midrule
离散质量函数 & $p(x)=P(X=x)$，各可能取值的质量相加为1。\\
概率密度 & 若分布具有密度$f$，则$F(t)=\int_{-\infty}^t f(u)\,du$；密度值不是单点概率。\\
累积分布函数 & 一般有$P(a<X\leq b)=F(b)-F(a)$，$a<b$。\\
\bottomrule
\end{tabularx}
\end{center}

记左极限为$F(a^-)=\lim_{t\uparrow a}F(t)$，则
\[
 P(X=a)=F(a)-F(a^-),\qquad
 P(a\leq X\leq b)=F(b)-F(a^-).
\]
例如左图$P(0<X\leq1)=1/4$，但$P(0\leq X\leq1)=1$；端点是否纳入会改变答案。
有密度的分布在单点上概率为0，此时有限个端点的开闭不改变区间概率。

\begin{pitfall}[几乎处处不等于处处]
把均匀分布密度在单个位置的值改成100，不改变任何区间积分，因而不改变分布。
若两个密度只在零测集（例如有限个点）上不同，便称为几乎处处（Almost Everywhere）相同。
类似地，ReLU只在0处不可微，不妨碍它在几乎处处可微；实现仍要约定0处如何反传。
\end{pitfall}

\clearpage
\section{标准误与一致性：更多样本能减少什么？}
\label{sec:math-standard-error}
标准差（Standard Deviation, SD）描述观测或随机变量自身的离散程度；
标准误（Standard Error, SE）描述\textbf{估计量在重复抽样中的}离散程度。
对独立同分布、均值$\mu$、有限方差$\sigma^2$的$n$个观测，
\[
 \mathbb E\bar X=\mu,\qquad
 \operatorname{Var}(\bar X)=\frac{\sigma^2}{n},\qquad
 \mathrm{SE}(\bar X)=\frac{\sigma}{\sqrt n}.
\]
$\sigma$未知时常用$s/\sqrt n$估计，$s^2=\sum_i(x_i-\bar x)^2/(n-1)$。
例如$[2,4,4,6]$的均值为4，$s^2=8/3$，$s\approx1.633$，均值的估计标准误约0.816。

\begin{figure}[H]
\centering\begin{tikzpicture}
\begin{axis}[width=135mm,height=44mm,xmin=-2,xmax=10,ymin=0,ymax=.9,
 xlabel={重复抽样得到的均值},ylabel={密度},xtick={0,2,4,6,8},
 legend style={font=\small,draw=none,at={(.98,.95)},anchor=north east}]
\addplot[Quiet,dashed,thick,domain=-2:10,samples=151] {exp(-(x-4)^2/8)/(2*sqrt(2*pi))};\addlegendentry{$n=1$，SE=2}
\addplot[Amber,thick,domain=-2:10,samples=151] {exp(-(x-4)^2/2)/sqrt(2*pi)};\addlegendentry{$n=4$，SE=1}
\addplot[Teal,very thick,domain=-2:10,samples=151] {exp(-(x-4)^2/.5)/(.5*sqrt(2*pi))};\addlegendentry{$n=16$，SE=0.5}
\end{axis}
\end{tikzpicture}

\caption{假设独立观测服从$N(4,4)$，则样本均值精确服从$N(4,4/n)$。曲线是不同样本量下均值的抽样密度。}
\end{figure}
\readfigure{$n=1,4,16$时均值的标准误为2、1、0.5。
分布收窄的是估计均值，不是单个新观测；样本量乘4，标准误才减半。}

\subsection{一致性不是“预测误差最终变成0”}
估计量$\hat\theta_n$依概率趋近真实参数$\theta$，称为一致：
对任意$\varepsilon>0$，$P(|\hat\theta_n-\theta|>\varepsilon)\to0$。
在上述条件下，切比雪夫不等式给出
\[
 P(|\bar X-\mu|>\varepsilon)\leq\frac{\sigma^2}{n\varepsilon^2}\longrightarrow0.
\]
所以样本均值一致。总体非退化时，只用第一个观测$X_1$估计$\mu$虽然无偏，却不会随$n$增加而稳定；
反之，$n\bar X/(n+1)$有限样本下通常有偏，仍可一致。

\begin{pitfall}[减少估计误差有条件]
相关或重复记录不能直接按独立样本数套用$1/\sqrt n$；分布变化也不由增加同类数据自动解决。
即使模型参数估计一致，单个新目标仍可能含不可约噪声，预测误差不必趋于0。
标准误本身也不是完整置信区间；区间构造需要分布或近似条件。
\end{pitfall}


\noindent\textbf{参考：}Deisenroth、Faisal与Ong，\textit{Mathematics for Machine Learning}（2020），
第2、4、5、6章。本附录图与数字例子独立设计；更多计算细节见本项目复核脚本。
